\chapter{Conclusion et Perspectives}

Ce projet livre une chaine complete de generation de narratives analogues,
depuis des textes Wikipedia bruts jusqu'a une application interactive deployable.
La branche \texttt{merge -> final\_en} a ete implementee (249 entites), mais les
resultats quantitatifs de similarite retenus dans ce rapport proviennent du
pipeline \texttt{en <-> fr\_en}.

La comparaison experimentale montre:
\begin{itemize}[leftmargin=2em]
  \item un gain net du fine-tuning face a la baseline non fine-tunee,
  \item une robustesse forte des modeles fine-tunes sur hard negatives,
  \item un compromis modele/seuil a expliciter selon le mode de deploiement.
\end{itemize}

Recommandation pratique:
\begin{itemize}[leftmargin=2em]
  \item usage recommande (meilleur score global): MPNet fine-tune avec seuil calibre (\(\approx 0.46\)),
  \item alternative plus legere: allMiniLM fine-tune avec seuil recommande (\(\approx 0.28\)).
\end{itemize}

Travaux futurs:
\begin{itemize}[leftmargin=2em]
  \item calibrage automatique du seuil selon la distribution de scores,
  \item extension a d'autres domaines (films, livres),
  \item evaluation humaine plus large de la qualite narrative.
\end{itemize}
